% !TeX spellcheck = pt_BR

\documentclass{abntex2}
\usepackage{abntex2cite}

\usepackage[brazil]{babel}
%sudo apt-get install texlive-lang-portuguese

\usepackage[T1]{fontenc} % caracteres latinos, turcos, etc

\usepackage{graphicx}
\graphicspath{ {.} }

\usepackage{color}
\usepackage{url}
\usepackage{epsfig}

\usepackage{algorithm,algpseudocode}

\newcommand{\specialcell}[2][c]{%
	\begin{tabular}[#1]{@{}c@{}}#2\end{tabular}}

\newcommand{\Comentar}[1]{\State {\cmmt{#1}}}
\newcommand{\Break}{\State \textbf{break}}
\renewcommand{\Return}{\State \textbf{return}~}
\renewcommand{\algorithmicensure}{\textbf{Parameters:}}

\usepackage{longtable}
\newcommand{\elite}{\mathcal{E}}
\usepackage{rotating}
%\usepackage{subfigure}
\usepackage{lscape}
\usepackage{multirow}
\usepackage{adjustbox}
\usepackage{comment}
\usepackage{setspace}
\usepackage{caption}
\captionsetup{labelfont=bf}
\captionsetup{skip=0pt} % vertical distance from the table or figure
\usepackage[margin=0.9in]{geometry}
\usepackage[table]{xcolor}
\usepackage{booktabs}

%\usepackage{natbib}


\begin{document}
		
	\title{Pesquisa Operacional em Suportabilidade}
	\author{Antonio Celio Pereira de Mesquita\\ AeroLogLab-ITA\\  \small{celio.mesquita@gp.ita.br} }
	\date{} % deixado vazio de propósito a fim de excluir a data

\maketitle

\begin{sloppypar}

	Este capítulo aborda os problemas inerentes à suportabilidade de sistemas, esboça modelos matemáticos e as possibilidades do uso das ferramentas de Pesquisa Operacional na sua solução.
	

\section{Introdução}

A capacidade de suporte de sistemas complexos refere-se à aptidão de um sistema tecnológico ou uma organização, de ser efetivamente mantido, atualizado, reparado e aprimorado temporalmente. Um sistema complexo é aquele sendo composto por muitas partes interconectadas que interagem de forma não linear, dificultando prever como o sistema se comportará em diferentes condições. Porém, um sistema complexo pode ser subdividido, para fins de análise, em subsistemas modeláveis e otimizáveis.

A capacidade de suporte envolve vários aspectos, incluindo a aptidão de diagnosticar problemas, a disponibilidade dos recursos necessários (como pessoal, ferramentas e peças de reposição), a eficiência dos processos de manutenção e reparo, a propriedade de fazer alterações no sistema sem causar interrupções e também de se adaptar a requisitos e ambientes em mudança.

Também é importante destacar que cada subsistema está exposto a riscos que precisam ser quantificados sob os pontos de vista da probabilidade de ocorrência de algum infortúnio, bem como do impacto que isto causaria ao sistema.

Um sistema complexo e seus subsistemas devem ser projetados com a capacidade de suporte em mente desde a concepção. Isso pode envolver o projeto do sistema com componentes modulares que podem ser facilmente substituídos ou atualizados, utilizando-se de ferramentas de diagnóstico eficazes, garantindo que os recursos necessários estejam disponíveis e projetando processos de manutenção e reparo que sejam eficientes. Além disso, o sistema deve ser continuamente monitorado e avaliado para identificar possíveis problemas, riscos e possibilidades de melhoria.

Melhores detalhes poderão ser encontrados em \cite{MarquesAbrahaoGuilherme:2020}.



\subsection{Aplicações da Pesquisa Operacional no desenho da suportabilidade}\

A pesquisa operacional (PO) é um campo de estudo que usa modelagem matemática, técnicas de otimização e métodos analíticos para resolver problemas complexos em negócios, engenharia e outros campos. O objetivo da PO é encontrar soluções ótimas ou muito próximas do ótimo em curto tempo para problemas que envolvem múltiplas variáveis, restrições e objetivos.

A PO pode ser usada para mitigar problemas de capacidade de suporte em sistemas complexos, fornecendo uma estrutura para analisar e melhorar a capacidade de suporte do sistema. Especificamente, a PO pode ser usada para:

\begin{itemize}
	
\item Prever a probabilidade e o impacto dos problemas de suporte: ao criar modelos que simulam o comportamento do sistema e identificam possíveis gargalos ou pontos de falha, a PO pode ajudar a prever onde os problemas de suporte são prováveis de ocorrer e qual seria o impacto dessas questões.

\item Otimizar a alocação de recursos: a PO pode ajudar a otimizar a alocação de recursos (como pessoal, ferramentas e peças sobressalentes) para garantir que os problemas de capacidade de suporte sejam tratados de forma eficiente.

\item Projetar processos eficazes de manutenção e reparo: a PO pode ser usada para projetar processos de manutenção e reparo que sejam eficientes e minimizem a interrupção do sistema.

\item Avaliar o impacto das mudanças: A PO pode ser usada para avaliar o impacto das mudanças propostas no sistema (como atualizações ou modificações) na capacidade de suporte e para identificar possíveis problemas antes que eles ocorram.

\end{itemize}

No geral, a PO fornece um poderoso conjunto de ferramentas e técnicas para mitigar problemas de capacidade de suporte em sistemas complexos e pode ajudar a garantir que esses sistemas permaneçam confiáveis.


\section{Organização científica da manutenção}


\subsection{O problema da programação da manutenção preventiva}

O problema de programação da manutenção preventiva (PPMP) pode ser definido como o problema relacionado com a programação da manutenção vivenciado por um sistema que possua uma frota de veículos (aviões, caminhões, maquinas, tratores, trens, etc.) sujeita a múltiplas restrições. O PPMP envolve dois tipos básicos de programação de atividades:

\begin{itemize}
	
	\item[I] Programação das tarefas de manutenção nas oficinas, onde os problemas são similares aos problemas de sequenciamento de tarefas (do inglês \textit{scheduling}), que envolvem, tipicamente, a alocação de recursos limitados a tarefas ao longo do tempo. Neste caso, a solução do problema está em encontrar uma sequência para todas as fases necessárias a uma revisão (testes, inspeções, delineamento das atividades, desmontagem, limpeza, reparos, substituições, montagem, pintura, testes, etc.), dadas limitações de recursos das oficinas.
	
	\item[II] Programação dos recolhimentos das aeronaves para as atividades de manutenção preventiva necessárias, segundo as restrições verificadas no sistema.
	
\end{itemize}

A solução, neste segundo caso, pouco ou nada tem a ver com a programação do que acontece nas oficinas, mas sim, com a ordenação dos recolhimentos para as atividades de manutenção preventiva necessários a uma frota de aeronaves ao longo de um período previsto de seu ciclo de vida. Este foi o foco do trabalho de  \cite{Abrahao:2006}.

\subsection{Decisões sobre internalizar ou não a manutenção de 3° nível}\

O processo de tomada de decisão sobre internalizar ou não a manutenção de 3º nível normalmente envolve a ponderação dos custos e benefícios associados a ambas as opções. Internalizar a manutenção pode ser uma opção mais econômica, pois elimina a necessidade de contratar e pagar pessoal de manutenção externo. Por outro lado, pode ser um processo mais demorado e complexo, pois requer treinamento e gerenciamento da equipe para garantir a qualidade do trabalho. Além disso, pode ser difícil obter peças e componentes com o mesmo custo de fornecedores de manutenção externos. No final das contas, a decisão deve ser baseada em uma modelagem e otimização dos custos, benefícios e riscos associados.
Normalmente, esta análise é realizada pelo AeroLogLab (\url{www.aerologlab.ita.br}) por meio da plataforma OPUS 10 (\url{www.systecongroup.com}).

\subsection{Análise de dados para diagnóstico ou prognóstico}\

A análise de dados para diagnóstico ou prognóstico para planejamento de manutenção de componentes de aeronaves normalmente envolve a coleta e análise de uma variedade de fontes de dados, como registros de voo e manutenção, previsões meteorológicas, aeronaves, sistemas e outras fontes de informações relacionadas à aeronave. Esses dados são usados para identificar possíveis problemas que podem afetar o desempenho e a segurança da aeronave, bem como para avaliar a condição de componentes e sistemas.

Essa análise é frequentemente usada para antecipar e planejar futuras tarefas de manutenção, como manutenção preventiva, reparos e substituições. Além disso, a análise de dados pode ser usada para monitorar o desempenho de aeronaves e identificar possíveis tendências ou padrões na manutenção de aeronaves.

\begin{figure}[h]
	\centering
	\includegraphics[width=0.5\textwidth]{Images/analise_dados.jpg}\\
	\small\textsuperscript{Origem: livre de direitos autorais em www.pexels.com}
\end{figure}

Informações relacionadas poderão ser encontradas em \cite{DuncanMarquesGuilhermeMartinetti:2005}.

\subsection{Avaliação do impacto de atrasos e cancelamentos em rotas operacionais}\

Os atrasos e cancelamentos em rotas operacionais de aeronaves comerciais envolve a análise de diversos fatores, como o custo direto e indireto, o tempo perdido para os passageiros, a satisfação dos clientes, a segurança e a eficiência operacional.

O custo direto inclui prejuízos financeiros como a perda de receita, custos de combustível, compensações aos passageiros, aumento de horas extras para os tripulantes e outras despesas. O custo indireto inclui os efeitos sobre a imagem da companhia aérea, a fidelidade dos passageiros e a reputação da companhia aérea. Além disso, a análise do impacto também pode incluir fatores como o número de atrasos e cancelamentos, a duração dos atrasos e a categoria de aeronave.

\begin{figure}[h]
	\centering
	\includegraphics[width=0.5\textwidth]{Images/flight_delay.jpg}\\
	\small\textsuperscript{Origem: livre de direitos autorais em www.pixabay.com}
\end{figure}

Um modelo matemático e uma ferramenta de PO pertinente podem auxiliar no planejamento, levando em conta os riscos de atraso ou cancelamento, com o objetivo de antever esses riscos e facilitar o planejamento.

\subsection{Plano de manutenção resiliente com otimização e aprendizado de máquina}\

\textit{Machine Learning} (ML) – ou aprendizado da máquina – está ligado às máquinas que possuem capacidade de aprendizagem própria para realizar ações ou tomar decisões para atividades através de técnicas específicas, \textit{Big Data} e inteligência artificial (\textit{IA}), \cite{Mitchell:1997}.  

O aprendizado pode ocorrer pela supervisão de um programador que faz a inserção de parâmetros no sistema, ou de forma interativa, onde os algoritmos adaptam-se aos dados, permitindo cruzar informações que façam previsões baseadas na análise de dados, sem terem sido programadas especificamente para isso.

Um plano de manutenção resiliente com otimização e aprendizado de máquina é uma abordagem proativa que usa análise preditiva e ciência de dados a fim de identificar possíveis problemas e otimizar os cronogramas de manutenção.

Essa categoria de plano usa algoritmos de aprendizado de máquina para analisar leituras de sensores, dados históricos e outras informações para prever quando o equipamento provavelmente precisará de manutenção e otimizar o tempo das tarefas de manutenção.

Essa abordagem pode melhorar a eficiência e a confiabilidade, além de reduzir o tempo de inatividade e os custos. Além disso, ajuda a identificar possíveis problemas antes que se tornem grandes, evitando assim reparos caros ou paradas não planejadas.

Conhecido como um problema de programação ou \textit{scheduling}, tarefas são agrupadas em pacotes temporais, algumas tarefas eventualmente poderão ser colocadas como \textit{out-of-phase} (fora dos pacotes) e, dentro de cada pacote, as tarefas são agrupadas por zona, por qualificação do mecânico e por precedência entre tarefas.
O objetivo global é a eficiência traduzida como minimização de custos e do tempo de parada, bem como minimização dos custos de oportunidade.

a Equação \ref{eq:prev} estabelece o custo da manutenção preventiva.

\begin{equation}\label{eq:prev}
prev_j =  pmc_j + pmoc_j,\ for\ j \in \{1, 2, 3, \ldots, m\}
\end{equation}


A Equação \ref{eq:corr} estabelece o custo da manutenção corretiva.

\begin{equation}\label{eq:corr}
corr_j =  cmc_j + cmoc_j,\ for\ j \in \{1, 2, 3, \ldots, m\}
\end{equation}


A Equação  \ref{eq:dist} estabelece a distância em termos temporais entre o limite de horas de voo e o momento da parada da aeronave.

\begin{equation}\label{eq:dist}
dist_{ij} = \lceil{\frac{stop_i}{max_j}}\rceil{}-\frac{stop_i}{max_j} 
\end{equation}


A Equação  \ref{eq:prob} determina a probabilidade de falha.

\begin{equation}\label{eq:prob}
prob_{ij} = 1 - e^{(-\lambda{}_j * stop_i)}
\end{equation}

As tarefas de preparação não podem ser duplicadas. As tarefas executam preparações comuns comente uma vez, a fim de não desperdiçar recursos.
$K_i$ é o conjunto de preparações do pacote $i$.

\begin{equation} \label{eq:min}
minimize\ =\ \sum_{i=1}^{m} \sum_{j=1}^{n} X_{ij} \times ( prev_j + dist_{ij}*prev_j + prob_{ij}*corr_j )
\end{equation}

A Equação \ref{eq:min} estabelece a Função Objetivo, que visa minimizar custos.


\subsection{Problema de planejamento de voos e manutenção}\

Do inglês \textit{flight maintenance planning problem} (FMPP), podem planejados pelas companhias aéreas por meio da divisão em grupos de problemas elementares: desenvolvimento das rotas, desenho do cronograma, designação de frota, rotas de manutenção e Pareamento de tripulação. 

As abordagens usadas para resolver esses sub-problemas podem ser divididas em dois grupos básicos: abordagens usadas apenas para resolver um dos problemas mencionados (abordagens de propósito único) e procedimentos que resolvem mais de um problema básicos em simultâneo (abordagens integradas).

O Problema de Planejamento de Voos e Manutenção é uma abordagem integrada de otimização que visa encontrar a melhor rota para o voo de um avião e a manutenção necessária para mantê-lo em condições seguras. Ele envolve determinar a rota ótima para o voo, bem como o calendário de manutenção apropriado para a aeronave antes de cada voo.

Está relacionado ao problema de programação de voos, mas tem um foco maior na manutenção e na segurança operacional. É importante para garantir que os voos sejam executados com segurança e eficiência, além de reduzir custos de manutenção. Este problema envolve as seguintes variáveis:

\vspace{0.5cm}

\begin{minipage}[h]{.98\linewidth}
	\begin{minipage}[t]{0.25\linewidth}
		\begin{tabular}{l}
		Tempo de voo\\
		Custo de combustível\\
		Consumo por hora de voo\\
		Tempo de manutenção\\
		Custo de manutenção\\
		Custo de oportunidade\\
		Restrições de segurança\\
		Restrições de custo\\
		Restrições de tempo\\
		Restrições de rota\\
		\end{tabular}
	\end{minipage}
	\hfill
	\begin{minipage}[t]{0.72\linewidth}
		\vspace{-2ex}
		\includegraphics[width=0.7\linewidth]{Images/programacao_voo.jpg}\\
		\small\textsuperscript{Origem: livre de direitos autorais em www.pexels.com}
	\end{minipage}
\end{minipage}



\subsection{Problema de roteamento e manutenção de aeronaves para frotas compartilhadas, incluindo informações de prognóstico de saúde das aeronaves}\

Os problemas de roteamento e manutenção de aeronaves são considerações críticas para frotas compartilhadas de aeronaves. Esses problemas afetam a segurança e a confiabilidade da aeronave e podem ter consequências financeiras significativas para a companhia aérea ou empresa que possui e opera a frota.

O roteamento de aeronave refere-se ao processo de determinação da rota ideal para uma aeronave viajar, considerando, fatores como consumo de combustível, condições climáticas e congestionamento do tráfego aéreo. Esse processo envolve a análise de abundantes de dados, incluindo dados históricos de voos, previsões meteorológicas e informações de tráfego aéreo em tempo real. O objetivo do roteamento de aeronaves é minimizar o tempo de voo, o consumo de combustível e os custos operacionais gerais, garantindo segurança e confiabilidade.

Uma maneira de resolver problemas de roteamento e manutenção de aeronaves é por meio do uso de sistemas de prognóstico de integridade de aeronaves. Esses sistemas usam sensores e análise de dados para monitorar a integridade e o desempenho dos componentes da aeronave em tempo real, permitindo que as equipes de manutenção detectem e resolvam possíveis problemas antes que se tornem problemas sérios. Os sistemas de manutenção preditiva também podem usar algoritmos de aprendizado de máquina para analisar dados históricos e prever quando a manutenção é necessária, ajudando a minimizar o tempo de inatividade e reduzir os custos.

Em resumo, os problemas de roteamento e manutenção de aeronaves são considerações críticas para frotas compartilhadas de aeronaves, e o gerenciamento eficaz dessas questões é essencial para garantir a segurança, confiabilidade e lucratividade da frota. O uso de sistemas de prognóstico de integridade de aeronaves pode ajudar a resolver esses problemas, fornecendo monitoramento em tempo real e recursos de manutenção preditiva.

Para detalhes sobre este assunto consulte \cite{Eduardo:2022a}.



\subsection{O problema da otimização do leiaute de uma instalação para manutenção de aeronaves}\

A otimização do leiaute de uma instalação para manutenção de aeronaves é uma parte importante do processo geral de manutenção de aeronaves. Este processo envolve a otimização do espaço na instalação para garantir que todos os equipamentos e pessoal necessários estejam localizados nas posições mais eficientes para maximizar a produtividade e minimizar os custos.

Isso pode envolver reorganizar o leiaute da instalação, como mover estações de trabalho e equipamentos, adicionar ou remover paredes e reconfigurar o fluxo de pessoal e materiais. Além disso, as instalações devem ser projetadas para reduzir riscos, como perigos elétricos e de incêndio, além de permitir acesso fácil a ferramentas e peças.

O leiaute ideal de uma instalação para manutenção de aeronaves também inclui considerar a ergonomia e a segurança do pessoal, fornecendo espaço adequado para movimentação e equipamentos de segurança.

Do ponto de vista da otimização, o objetivo seria minimizar os percursos de movimentação dos técnicos, conferindo agilidade e eficiência. As ferramentas mais utilizadas tenderiam a ficar mais acessíveis e as pouco empregadas em posições mais recuadas.

O modelo matemático geralmente utilizado é o de uma árvore geradora mínima (\textit{ minimum spanning tree}).

Dada um grafo (G, $\phi$), pode-se definir o peso de um subgráfo $H \subset G$ como:

\begin{equation} \label{eq:fo1}
\phi(H) \gets \sum_{e \in E(H)} \phi (e)
\end{equation}

$E(H)$ é o subconjunto de arestas $e$ de $H$.

Em um grafo conectado $G$, uma árvore geradora mínima $T$ é uma árvore com valor mínimo.

\begin{equation} \label{eq:fo2}
min_T\ H(T) \gets \sum_{e \in E(H)} \phi (e) \times X_{ij}
\end{equation}


Considere $X_{ij}$ igual a 1 se a aresta (i,j) está na árvore, 0 caso contrário.


\subsection{Otimização da localização de instalações de manutenção}\

A otimização da localização de instalações de manutenção de aeronaves é a análise de fatores para determinar o melhor local para a instalação de um hangar de manutenção.

Estes fatores incluem a proximidade de outras instalações de manutenção, a quantidade de tráfego aéreo, as condições climáticas, a infraestrutura disponível, as qualificações de mecânicos disponíveis e outros fatores. Além disso, existem várias variáveis que devem ser consideradas, como custos, localização e segurança.

O objetivo é maximizar a eficiência global da companhia aérea, minimizando custos.

Um modelo simplificado pode ser semelhante a este:


\begin{equation} \label{eq:fo3}
minimize \sum_{i \in M} \sum_{j \in N} c_{ij} \times X_{ij} + \sum_{i \in M} f_i \times z_i
\end{equation}

onde:

$X_{ij}$ é a fração da demanda do cliente $j$ fornecida pela instalação $i$

$z_i$ = 1 se a instalação estiver aberta, 0 caso contrário

$c_{ij}$ é o custo do serviço e distribuição para fornecer aos clientes a partir da instalação $i$

$f_i$ é o custo fixo de construir a instalação $i$

$M$, $N$ são os conjuntos de possíveis locais para as instalações e a localização dos clientes.


\section{Organização científica do suprimento}


\subsection{Otimização de estoques de peças sobressalentes}\

A otimização de peças de reposição é um processo utilizado para minimizar a quantidade de peças de reposição que uma empresa precisa ter em seu estoque a fim de atender à demanda dos clientes, especialmente, no nosso caso, as oficinas de manutenção. Envolve analisar os vários tipos de peças estocadas e determinar quais peças são realmente necessárias e quais possuem baixa taxa de consumo. Eventualmente, compensa ser reativo e adquirir somente quando surge a necessidade. A não ser quando é sabido que determinados fornecimentos são demorados e o componente precisa ser mantido em estoque a fim de evitar disrupção nas operações.

A otimização do estoque de peças sobressalentes pode ajudar uma empresa a economizar dinheiro em custos de armazenamento, reduzir o desperdício e melhorar o atendimento ao cliente, tendo as peças necessárias prontamente disponíveis quando os clientes precisarem delas, seja oriunda do estoque ou de um fornecedor local.

Assumindo que a demanda é conhecida a partir dos dados do ano anterior e que seria constante para uma temporada inteira, mas diferente para diferentes temporadas.

Um modelo simples que representa essa situação é o da quantidade de pedido econômico (\textit{ economic order quantity} - EOQ), às vezes também é chamado de modelo econômico de tamanho de lote.

O modelo EOQ é a quantidade do pedido que minimiza os custos totais de manutenção e os custos de pedido.

Exemplo de Quantidade de Pedido Econômico:

O EOQ considera o momento do novo pedido, o custo incorrido para fazer um pedido e os custos para armazenar mercadorias. Se uma empresa está constantemente fazendo pequenos pedidos para manter um nível de estoque específico, os custos de pedidos são maiores, com a necessidade de espaço de armazenamento adicional.

\begin{equation} \label{eq:eoq}
	EOQ \gets \sqrt{\frac{2 \times S \times D}{H}}
\end{equation}

Onde $S$ é o custo por pedido, $D$ a quantidade vendida por ano e $H$ o custo de manutenção do estoque.

Por exemplo, considere um fornecedor de componentes eletrônicos que oferece uma linha de fones aeronáuticos. Ele vende 1.000 fones por ano. Custa à empresa \$5 por ano para manter um único fone em estoque, e o custo fixo para fazer um pedido é de \$2.

A fórmula EOQ é a raiz quadrada de (2 x 1.000 fones x custo do pedido de \$2) / (custo de manutenção de \$5), ou 28,3 com arredondamento. O tamanho ideal do pedido para minimizar custos e atender a demanda dos clientes é de pouco mais de 28 fones.

\subsection{Estimativas de abastecimento por meio de aprendizado de máquina}\

\textit{Machine Learning} (ML) – ou aprendizado da máquina – está ligado às máquinas que possuem capacidade de aprendizagem própria para realizar ações ou tomar decisões para atividades através de técnicas específicas, \textit{Big Data} e inteligência artificial (\textit{IA}), \cite{Mitchell:1997}.

O aprendizado pode ocorrer pela supervisão de um programador que faz a inserção de parâmetros no sistema, ou interativamente, no qual os algoritmos adaptam-se aos dados, permitindo cruzar informações que façam previsões baseadas na análise de dados, sem que tenham sido programadas especificamente para isso.

\textbf{Importância}: A importância do ML no controle de estoque é projetar a demanda com maior segurança e precisão, prevendo fornecimentos futuros sem grandes riscos e viabilizando pedidos de compras automaticamente. Além disso, o ML pode prever perdas e solucionar problemas que comprometem os resultados, como rupturas, custos elevados com carregamentos e baixo giro dos produtos em suas respectivas curvas.

\textbf{Vantagens}: A aprendizagem de máquina influencia as atividades mais complexas da logística e dos armazéns, proporcionando maior precisão e transparência nos dados.

\textbf{Desempenho da entrega}: A eficiência na gestão de entregas e os custos de frete podem ser aperfeiçoados com a utilização do ML. Com parâmetros e uma extensa base de dados, a escolha dos métodos de transporte é mais precisa, diminuindo os riscos para toda a cadeia de suprimentos.

\textbf{Previsão da demanda}: Com o aprendizado de máquina, a previsão da demanda se torna mais definida, auxiliando no planejamento estratégico de vendas e logística e permitindo que as operações trabalhem com baixos níveis de estoque.

\textbf{Visibilidade de toda a cadeia de suprimento}: O ML, juntamente à telemática (Internet das Coisas - loT), permite que toda a logística seja monitorada em tempo real, aumentando os resultados das entregas e a comunicação entre setores.

\subsection{Simulação dos processos internos de suprimento e otimização de leiaute}\

Os processos internos e a otimização de leiaute pode ser simulada e testada para determinar o melhor leiaute possível para um determinado espaço.

A simulação de processos internos e otimização de leiaute podem ajudar a reduzir custos, melhorar o atendimento ao cliente do armazém e aumentar a eficiência da cadeia de suprimentos.

Do ponto de vista da otimização, o objetivo seria minimizar os percursos de movimentação dos suprimentistas, conferindo agilidade e eficiência. Os componentes mais fornecidos tenderiam a ficar mais acessíveis e os pouco fornecidos em armários mais recuados e prateleiras mais altas.

O modelo matemático geralmente utilizado é o de uma árvore geradora mínima (\textit{minimum spanning tree}).

Dada um grafo (G, $\phi$), pode-se definir o peso de um subgrafo $H \subset G$ como:

\begin{equation} \label{eq:fo4}
	\phi(H) \gets \sum_{e \in E(H)} \phi (e)
\end{equation}

$E(H)$ é o subconjunto de arestas $e$ de $H$.

Em um grafo conectado $G$, uma árvore geradora mínima $T$ é um grafo com valor mínimo.

\begin{equation} \label{eq:fo5}
	min_T\ H(T) \gets \sum_{e \in E(H)} \phi (e) \times X_{ij}
\end{equation}


Considere $X_{ij}$ igual a 1 se a aresta (i,j) está na árvore, 0 caso contrário. As arestas são os percursos entre o ponto de entrega das mercadorias e os armários ou prateleiras.



\subsection{Otimização da localização de instalações de distribuição}\

A otimização da localização das instalações de distribuição é um aspecto importante do gerenciamento da cadeia de suprimentos. Ao localizar estrategicamente os centros de distribuição próximos ao ponto de demanda, as empresas podem reduzir os custos de transporte, minimizar os prazos de entrega e melhorar o atendimento ao cliente.

A otimização da localização das instalações de distribuição geralmente é feita usando uma combinação de modelos analíticos, como sistemas de informações geográficas, para identificar os locais mais adequados e, em seguida, usar algoritmos de otimização para escolher a localização ideal. Esse processo pode ajudar a garantir que o centro de distribuição esteja localizado no local mais eficiente e econômico.

Localizações não otimizadas pode acarretar ineficiência na suportabilidade de sistemas complexos.

\subsection{Análise multicritério na seleção de fornecedores}\

As análises multicritério geralmente utilizam os métodos \textit{analytic network process} (ANP) e \textit{analytic hierarchy process} (AHP).

O ANP é uma forma mais geral do AHP usado na análise de decisão multicritério.

O AHP estrutura um problema de decisão em uma hierarquia com um objetivo, critérios de decisão e alternativas, enquanto o ANP o estrutura como uma rede. Ambos então usam um sistema de comparações pareadas para medir os pesos dos componentes da estrutura e, finalmente, classificar as alternativas na decisão.

Por exemplo:

A fim de se contratar um engenheiro, pode-se elencar duas competências: experiência e habilidades em gestão.

A experiência pode ser técnica ou de vendas, e as habilidades de gestão podem ser: da produção, de equipes, de finanças e de manutenção.

Supondo-se que haja três candidatos, compara-se pares de competências, pontuando-se de 0 a 9 para um candidato ou para o outro, sendo que, no final, uma operação entre matrizes define graus de acordo com todas as avaliações par-a-par realizadas por um ou por mais de um avaliador.

Um exemplo prático pode ser encontrado em \url{tool.aerologlab.ita.br/mcdm2}.

No caso de uma seleção de fornecedores, será preciso definir o objetivo, os critérios e subcritérios de avaliação, bem como os fornecedores candidatos.
Detalhes do funcionamento podem ser encontrados em \cite{Saaty:1996} e \cite{Saaty:2005}.


\subsection{Problema de roteirização como fator impactante da suportabilidade}\

O fornecimento de peças de reposição muitas vezes dá pouca atenção à otimização das rotas de transporte porque muitas vezes é visto como uma preocupação secundária e não como uma prioridade.
Além disso, os métodos existentes de otimização de rotas de transporte nem sempre são considerados do ponto de vista da capacidade de suporte. Como resultado, muitas organizações simplesmente optam pela rota mais direta, que muitas vezes não considera outros fatores, como recursos disponíveis, custo e tempo.

Adicionalmente, a complexidade do processo e a falta de recursos e competência também podem contribuir para a falta de atenção dada à suportabilidade na otimização das rotas de transporte.

Parte desse processo de distribuição pode ser terceirizado ou realizado por meio de frota própria, a depender da estratégia da organização.

De qualquer forma, faz-se mister que as demandas de cada cliente sejam atendidas, trecho a trecho, de possíveis rotas de coleta e entrega de componentes e equipamentos.

O problema, então, consiste em realizar rotas de coleta e entrega simultâneas com restrições de capacidade em volume,  peso e passeio do centro de gravidade do veículo (aeronave, caminhão ou navio). Além disso, restrições de destino, a fim de que componentes e equipamentos sejam entregues ao cliente correto. Este foi o foco do trabalho de doutorado de \cite{CelioMesquita:2023}.


\subsection{O problema da gestão de obsolescência}\

O gerenciamento de obsolescência de componentes envolve identificar e gerenciar os riscos associados a componentes e peças obsoletos ou em vias de faltarem fornecedores. Isso inclui monitorar os níveis de consumo e estoque de componentes e peças, prever a demanda futura, planejar fornecedores alternativos e desenvolver estratégias para substituir componentes obsoletos.

Destaca-se, ainda, que as organizações devam considerar o desenvolvimento de relacionamentos com fornecedores que ajudem a reduzir os riscos de obsolescência, seja por meio de engenharia reversa ou de modelos atualizados que superem as funcionalidades dos antigos.


\section{Desenvolvimento da suportabilidade}\


\subsection{Otimização e simulação da suportabilidade de um novo produto}\

\subsection{Análise e simulação da degradação por questões de suportabilidade}\

\subsection{problema de gerenciamento dos desempenhos de suportabilidade}\

\subsection{problema de alocação de aeronaves à luz do desempenho individualizado em termos de confiabilidade de cada aeronave}\

\subsection{problema de imaturidade nos fatores nos fatores RAMS}\

\subsection{Garantia da confiabilidade, da capacidade de manutenção e da capacidade de suporte sob logística baseada em desempenho}\

\subsection{Previsão das principais métricas de desempenho e custo para um novo sistema}\


%\bibliographystyle{elsarticle-harv}

\bibliographystyle{authordate1}

\bibliography{livro}

\end{sloppypar}

\end{document}
